\clearpage
\VerdierPhysicalPageAnchor{0214}%
\begingroup\small\setlength{\parskip}{0pt}\setlength{\emergencystretch}{1.5em}
\setlength{\abovedisplayskip}{4pt plus 2pt minus 2pt}%
\setlength{\belowdisplayskip}{4pt plus 2pt minus 2pt}%

\VerdierSemanticAnchor{ch3d-3-2-5-b}
\noindent b) 次の図において：

\VerdierSemanticAnchor{ch3d-3-2-5-4}
\begin{equation}
\resizebox{0.96\textwidth}{!}{%
$\begin{tikzcd}[ampersand replacement=\&,column sep=4.1em,row sep=3.0em]
\&\& Z''^\bullet
  \arrow[ddll,bend right=8,"{u(S^{m+n})}"']
  \arrow[dl,"\rho"']
  \arrow[dr,"\psi"]
  \arrow[ddrr,bend left=8,"{v(S^{m+n})}"]
\&\& \\
\& Z^\bullet
  \arrow[dl,"{u(S^n)}"']
  \arrow[dr,"{v(S^n)}"]
\&\& {Z'}^\bullet[n]
  \arrow[dl,"{u(S^m)[n]}"']
  \arrow[dr,"{v(S^m)[n]}"]
\& \\
X \&\& X'[n] \&\& X''[m+n]
\end{tikzcd}$%
}
\tag{3.2.5.4}
\end{equation}
中央の正方形と両側の三角形は $\operatorname{comp}(\mathcal A)$ で
可換である。

\VerdierSemanticAnchor{ch3d-3-2-5-proof-conclusion}
圏 $D(\mathcal A)$ に移ると、図 (3.2.5.4) から次の等式を得る：
\[
\bigl(\delta^m(S^m(X',X''))\bigr)[n]\,
\delta^n(S^n(X,X'))
=\delta^{m+n}(S^{m+n}(X,X'')).
\]
したがって (3.2.3.2) を用いれば証明が完了する。

\VerdierSemanticAnchor{ch3d-3-2-6}
\medskip\noindent\textbf{3.2.6.}\quad
命題 (3.2.5) の特に重要な応用は次のものである。次を
\[
S^n(X,Y)=0\longrightarrow Y\xrightarrow{u_n}Z_{n-1}
\xrightarrow{u_{n-1}}\cdots\xrightarrow{u_1}Z_0
\xrightarrow{u_0}X\longrightarrow0
\]
$E_{\mathcal A}^{n}(X,Y)$ の元とする。各整数 $p$，
$1\leq p\leq n$ に対して次のようにおく：

\VerdierSemanticAnchor{ch3d-3-2-6-1}
\begin{equation}
S_p^1=0\longrightarrow\operatorname{Ker}(u_{p-1})
\longrightarrow Z_{p-1}\longrightarrow\operatorname{Im}(u_{p-1})
\longrightarrow0.
\tag{3.2.6.1}
\end{equation}
このとき次の等式が成り立つ：

\VerdierSemanticAnchor{ch3d-3-2-6-2}
\begin{equation}
\delta^n(S^n(X,Y))=(-1)^{\frac{n(n-1)}2}
\delta^1(S_n^1)[n-1]\delta^1(S_{n-1}^1)[n-2]\cdots
\delta^1(S_1^1),
\tag{3.2.6.2}
\end{equation}
これは命題 (3.2.5) を用いれば直ちに従う。

\VerdierSemanticAnchor{ch3d-3-2-7}
\medskip\noindent\textbf{3.2.7.}\quad
ここで苦杯を最後まで飲み干そう。$X$ と $Y$ を $\mathcal A$ の二つの
対象、したがって $\mathcal A^o$、すなわち $\mathcal A$ の反対圏の
対象とする。集合 $E_{\mathcal A}^{n}(X,Y)$ と
$E_{\mathcal A^o}^{n}(Y,X)$ は等しく、二つの写像がある：

\VerdierSemanticAnchor{ch3d-3-2-7-1}
\begin{equation}
\resizebox{0.74\textwidth}{!}{%
$\begin{tikzcd}[ampersand replacement=\&,column sep=5.0em,row sep=2.8em]
E_{\mathcal A}^{n}(X,Y)
  \arrow[r,"{\delta_{\mathcal A}^{n}}"]
  \arrow[d,equal]
\& \operatorname{Ext}_{\mathcal A}^{n}(X,Y) \\
E_{\mathcal A^o}^{n}(Y,X)
  \arrow[r,"{\delta_{\mathcal A^o}^{n}}"']
\& \operatorname{Ext}_{\mathcal A^o}^{n}(Y,X)
\end{tikzcd}$%
}
\tag{3.2.7.1}
\end{equation}
\endgroup
